\documentclass[12pt]{article}
\usepackage{amsmath, amssymb}
\usepackage[margin=1in]{geometry}
\setlength{\parskip}{6pt}
\title{QUBO Formulation: Maximum Clique Problem}
\date{}
\begin{document}
\maketitle

\subsection*{Maximum Clique Problem}

\paragraph{Problem Statement.}
Given an undirected graph $G = (V, E)$ with vertex set $V$ and edge set $E$, a clique is a subset $C \subseteq V$ such that every pair of distinct vertices in $C$ is connected by an edge in $E$. The Maximum Clique Problem seeks to find a clique $C^*$ of maximum cardinality, i.e., $|C^*| = \max_{C \subseteq V, C \text{ is a clique}} |C|$.

\paragraph{Decision Variables.}
For each vertex $i \in V$, introduce a binary variable $x_i \in \{0, 1\}$. The variable $x_i$ takes the value $1$ if vertex $i$ is included in the candidate clique, and $0$ otherwise. The decision vector is $\mathbf{x} = (x_i)_{i \in V} \in \{0, 1\}^{|V|}$.

\paragraph{QUBO Derivation.}
Step 1 --- Write the original objective in general notation and flip the sign for minimisation. The problem requires maximising the number of selected vertices, which is equivalent to minimising its negative:
\begin{align}
  \min_{\mathbf{x} \in \{0, 1\}^{|V|}} \left( -\sum_{i \in V} x_i \right).
\end{align}

Step 2 --- Write each constraint in its original algebraic form using general indices. A subset of vertices forms a clique if and only if no two selected vertices are non-adjacent. For every pair of distinct vertices $(i, j)$ that are not connected by an edge, we enforce:
\begin{align}
  x_i + x_j \le 1, \quad \forall (i, j) \notin E, \ i < j.
\end{align}

Step 3 --- Convert each constraint to a penalty term. For binary variables, the inequality $x_i + x_j \le 1$ is violated if and only if $x_i = 1$ and $x_j = 1$. The product $x_i x_j$ equals $1$ exactly on violations and $0$ otherwise. We therefore encode the constraint directly as $\lambda x_i x_j$. We explicitly avoid the squared form $(x_i + x_j - 1)^2$, as it evaluates to $1$ for both $(1,1)$ and $(0,0)$, incorrectly penalizing valid assignments. The total penalty contribution is:
\begin{align}
  \lambda \sum_{(i,j) \notin E, \ i < j} x_i x_j.
\end{align}

Step 4 --- Combine objective and all penalty terms into the single QUBO Hamiltonian $H(\mathbf{x})$:
\begin{align}
  H(\mathbf{x}) = -\sum_{i \in V} x_i + \lambda \sum_{(i,j) \notin E, \ i < j} x_i x_j.
\end{align}

Step 5 --- Read off the Q matrix entries and offset $c$ from $H$ in general closed form. Matching $H(\mathbf{x}) = \sum_{i \in V} Q_{i,i} x_i + \sum_{i < j} Q_{i,j} x_i x_j + c$, we obtain:
\begin{align}
  Q_{i,i} &= -1, & \forall i \in V, \\
  Q_{i,j} &= \begin{cases} \lambda & \text{if } (i,j) \notin E, \\ 0 & \text{if } (i,j) \in E, \end{cases} & \forall i, j \in V, \ i < j, \\
  c &= 0.
\end{align}

\paragraph{Penalty Weight Justification.}
The objective function assigns a gain of $1$ for each vertex added to the clique. A single constraint violation (selecting two non-adjacent vertices) incurs a penalty of $\lambda$. To guarantee that the global minimum corresponds to a feasible clique, the penalty must strictly exceed the maximum possible objective benefit of violating a constraint. Since the maximum single-variable objective coefficient magnitude is $1$, setting $\lambda > 1$ ensures any violation outweighs the best possible objective gain from adding a vertex. We adopt $\lambda = 2$ as a robust integer choice that satisfies this condition for any instance size $n = |V|$.

\paragraph{Correctness Argument.}
Minimising $H(\mathbf{x})$ over $\{0, 1\}^{|V|}$ recovers the optimal solution because (i) any feasible clique assignment satisfies all constraints, yielding a penalty of zero and an energy equal to $-|C|$; (ii) any infeasible assignment violates at least one non-edge constraint, incurring a penalty of at least $\lambda > 1$. Since the maximum possible reduction in the objective term from adding an additional vertex is $1$, the penalty strictly dominates any infeasible configuration's potential energy advantage. Consequently, the global minimum of $H$ must be attained at a feasible assignment with the largest possible cardinality, i.e., the maximum clique.

\paragraph{Illustrative Example.}
Consider a graph with $n = 3$ vertices $\{0, 1, 2\}$ and no edges ($E = \emptyset$). All pairs are non-adjacent, so the constraint $x_i + x_j \le 1$ applies to $(0,1)$, $(0,2)$, and $(1,2)$. Using $\lambda = 2$, the general formulas instantiate to:
\begin{align}
  H(x_0, x_1, x_2) &= -x_0 - x_1 - x_2 + 2x_0 x_1 + 2x_0 x_2 + 2x_1 x_2.
\end{align}
The concrete Q matrix entries are $Q_{0,0}=Q_{1,1}=Q_{2,2}=-1$, $Q_{0,1}=Q_{0,2}=Q_{1,2}=2$, and $c=0$. Evaluating all $2^3 = 8$ binary assignments, the minimum energy is $-1$, achieved at $(1,0,0)$, $(0,1,0)$, and $(0,0,1)$. These correspond to cliques of size $1$, which is optimal for an edgeless graph.

\end{document}